% Source: Supplement: lower-bound circulation reduction and source/sink flow.
% Review and validation: docs/validation.md.
\subsection{上下界流：可行循环流与源汇流}
\label{math:bounded-flow}
\subsubsection*{有下界时怎样恢复守恒}
每条边要求 $l_e\leq f_e\leq r_e$，先固定下界流量，再把剩余容量设为 $r_e-l_e$。对每点维护
\[
    b_v=\sum_{e\text{ 入 }v}l_e-\sum_{e\text{ 出 }v}l_e.
\]
$b_v>0$ 表示下界已带来过多流入，剩余流应从该点流出。添加超级源 $SS$、超级汇 $TT$：$b_v>0$ 连 $SS\to v$，容量 $b_v$；$b_v<0$ 连 $v\to TT$，容量 $-b_v$。

原图存在可行循环流，当且仅当 $SS\to TT$ 最大流等于 $\sum_{b_v>0}b_v$，也即超级源的出边全部饱和。方向可通过检查守恒推导，不能把入下界与出下界的符号混用。原边最终流量为下界加上剩余网络在该边上的净流量。

\subsubsection*{使用教程：可行循环流示例}
先复制 Dinic。以下边表依次为起点、终点、下界、上界；所有下界非负且不超过上界，平衡量及所需总流量须在 64 位范围和 Dinic 上限内。保存原边编号，以便恢复具体流量。
\begin{lstlisting}[style=mtycpp,numbers=none]
int n=3,ss=n+1,tt=n+2;
Dinic g;
g.init(n+2);
vector<tuple<int,int,long long,long long>> e={
    {1,2,1,3},{2,3,0,2},{3,1,0,2}
};
vector<long long> b(n+1,0),f;
vector<int> id;
for(auto [u,v,l,r]:e){
    id.push_back(g.add(u,v,r-l));
    b[u]-=l,b[v]+=l;
}
long long need=0;
for(int u=1;u<=n;u++){
    if(b[u]>0){g.add(ss,u,b[u]);need+=b[u];}
    if(b[u]<0) g.add(u,tt,-b[u]);
}
bool ok=g.flow(ss,tt)==need;
if(ok){
    for(int i=0;i<(int)e.size();i++)
        f.push_back(get<2>(e[i])+g.e[id[i]^1].cap);
}
\end{lstlisting}

\subsubsection*{有源汇上下界流}
先额外加入 $t\to s$ 的大容量边，把源汇流闭合成循环流，再做上述可行性检查。成功时，该附加边上的流量记为 $F_0$。保存所有辅助边的编号，随后将\textbf{附加边与全部超级源汇边的两个残量方向都置零}，保留原图已有流量的残量状态。

在此基础上，从 $s$ 向 $t$ 继续增广得到最大可行流 $F_0+\Delta$；若要求非负的最小可行流，从 $t$ 向 $s$ 最多撤销 $F_0$，得到 $F_0-\Delta$。该构造以存在非负源汇流为前提；若题目允许负流值，需要重新处理辅助边与目标定义。

若要求恰好 $F$ 单位源汇流，可以把附加的 $t\to s$ 边上下界均设为 $F$，然后判定循环流可行性。带费用的上下界最优流还需处理循环流的费用优化，不能仅在可行性网络上随意调用普通最大流就得到最小费用。
